% GENERATED FILE. Edit canonical lessons, then run npm run generate:books.
% canonical-source-hash: fnv1a64:28818a46dfd58a2e
% canonical-lessons: FR-C411-percuter, FR-C411-heurter, FR-C411-pieton, FR-C411-cycliste, FR-C411-constater

\chapter{Un petit accident}
\label{ch:fr-un-petit-accident}
\hlchaptermodality{411}

\begin{chapteropening}
\textbf{By the end of this chapter:} \emph{I can describe a small road accident: hitting something, bumping into it, the pedestrian, the cyclist and writing up what happened.}

Complete the last lesson of chapter 411: I can describe a small road accident: hitting something, bumping into it, the pedestrian, the cyclist and writing up what happened.
\end{chapteropening}

\section[percuter]{percuter — to crash into, to hit}

\label{lesson:FR-C411-percuter}



\begin{quote}
Before the new one: say the French for the rest area, then the French for to overtake.
\end{quote}

\subsection*{You'll want to know: percuter}
\textbf{percuter} — \textquotedblleft{}to crash into, to hit\textquotedblright{}.

\begin{quote}
Une moto a percuté l'arrière du bus.
\end{quote}

\begin{quote}
\emph{A motorbike crashed into the back of the bus.}
\end{quote}

A hard, violent hit.

\begin{cousinweb}
\emph{percuter} is a learned borrowing from Latin \emph{percutere} 'to strike hard' (EtymDB), from \emph{per-} and \emph{quatere} 'to shake'.
\end{cousinweb}

\begin{grammarlens}[title={What you've built}]
One more word for telling, explaining and describing at B1.
\end{grammarlens}

\subsection*{Your turn}
\begin{itemize}
\raggedright
  \item \emph{Say it:} \emph{percuter}
  \item \emph{Say it:} \emph{percuter}, once more
  \item \emph{Say it:} a sentence of your own with \emph{percuter}
  \item \emph{Recall:} say the French for the rest area, then the French for to overtake, then say \emph{percuter} again
\end{itemize}

\begin{tcolorbox}[breakable,colback=teal!4,colframe=teal!35!black,title={Before you move on}]
What does percuter mean? (\textquotedblleft{}To crash into, to hit\textquotedblright{}.) Say it once more.
\end{tcolorbox}
\section[heurter]{heurter — to bump into, to knock against; to offend}

\label{lesson:FR-C411-heurter}



\begin{quote}
Before the new one: say the French for to overtake, then the French for to crash into.
\end{quote}

\subsection*{You'll want to know: heurter}
\textbf{heurter} — \textquotedblleft{}to bump into, to knock against; to offend\textquotedblright{}.

\begin{quote}
En reculant, il a heurté un poteau.
\end{quote}

\begin{quote}
\emph{While reversing he bumped into a post.}
\end{quote}

Also figurative: \emph{heurter quelqu'un} — to hurt someone's feelings.

\begin{cousinweb}
\emph{heurter} comes from Old French \emph{hurter} 'to ram', which is built on Frankish \emph{hūrt} 'battering ram' (EtymDB).
\end{cousinweb}

\begin{grammarlens}[title={What you've built}]
One more word for telling, explaining and describing at B1.
\end{grammarlens}

\subsection*{Your turn}
\begin{itemize}
\raggedright
  \item \emph{Say it:} \emph{heurter}
  \item \emph{Say it:} \emph{heurter}, once more
  \item \emph{Say it:} a sentence of your own with \emph{heurter}
  \item \emph{Recall:} say the French for to overtake, then the French for to crash into, then say \emph{heurter} again
\end{itemize}

\begin{tcolorbox}[breakable,colback=teal!4,colframe=teal!35!black,title={Before you move on}]
What does heurter mean? (\textquotedblleft{}To bump into, to knock against; to offend\textquotedblright{}.) Say it once more.
\end{tcolorbox}
\section[le piéton]{le piéton — the pedestrian}

\label{lesson:FR-C411-pieton}



\begin{quote}
Before the new one: say the French for to crash into, then the French for to bump into.
\end{quote}

\subsection*{You'll want to know: le piéton}
\textbf{le piéton} — \textquotedblleft{}the pedestrian\textquotedblright{}.

\begin{quote}
Les voitures s'arrêtent pour laisser passer les piétons.
\end{quote}

\begin{quote}
\emph{The cars stop to let the pedestrians cross.}
\end{quote}

\emph{Une rue piétonne} — a pedestrian street.

\begin{cousinweb}
\emph{piéton} is built on \emph{pied} 'foot' (EtymDB), which is inherited from Latin \emph{pēs} 'foot'.
\end{cousinweb}

\begin{grammarlens}[title={What you've built}]
One more word for telling, explaining and describing at B1.
\end{grammarlens}

\subsection*{Your turn}
\begin{itemize}
\raggedright
  \item \emph{Say it:} \emph{le piéton}
  \item \emph{Say it:} \emph{le piéton}, once more
  \item \emph{Say it:} a sentence of your own with \emph{le piéton}
  \item \emph{Recall:} say the French for to crash into, then the French for to bump into, then say \emph{le piéton} again
\end{itemize}

\begin{tcolorbox}[breakable,colback=teal!4,colframe=teal!35!black,title={Before you move on}]
What does le piéton mean? (\textquotedblleft{}The pedestrian\textquotedblright{}.) Say it once more.
\end{tcolorbox}
\section[le cycliste]{le cycliste — the cyclist}

\label{lesson:FR-C411-cycliste}



\begin{quote}
Before the new one: say the French for to bump into, then the French for the pedestrian.
\end{quote}

\subsection*{You'll want to know: le cycliste}
\textbf{le cycliste} — \textquotedblleft{}the cyclist\textquotedblright{}.

\begin{quote}
Le cycliste roulait sur la droite.
\end{quote}

\begin{quote}
\emph{The cyclist was riding on the right.}
\end{quote}

Same form for a woman: \emph{une cycliste}.

\begin{cousinweb}
\emph{cycliste} is built on \emph{cycle} 'bicycle', a learned word from Latin \emph{cyclus}, from Greek \emph{kýklos} 'circle, wheel' (Lewis \& Short).
\end{cousinweb}

\begin{grammarlens}[title={What you've built}]
One more word for telling, explaining and describing at B1.
\end{grammarlens}

\subsection*{Your turn}
\begin{itemize}
\raggedright
  \item \emph{Say it:} \emph{le cycliste}
  \item \emph{Say it:} \emph{le cycliste}, once more
  \item \emph{Say it:} a sentence of your own with \emph{le cycliste}
  \item \emph{Recall:} say the French for to bump into, then the French for the pedestrian, then say \emph{le cycliste} again
\end{itemize}

\begin{tcolorbox}[breakable,colback=teal!4,colframe=teal!35!black,title={Before you move on}]
What does le cycliste mean? (\textquotedblleft{}The cyclist\textquotedblright{}.) Say it once more.
\end{tcolorbox}
\section[constater]{constater — to note, to record (what happened)}

\label{lesson:FR-C411-constater}



\begin{quote}
Before the new one: say the French for the pedestrian, then the French for the cyclist.
\end{quote}

\subsection*{You'll want to know: constater}
\textbf{constater} — \textquotedblleft{}to note, to record (what happened)\textquotedblright{}.

\begin{quote}
Les deux conducteurs ont constaté les dégâts ensemble.
\end{quote}

\begin{quote}
\emph{The two drivers noted the damage together.}
\end{quote}

\emph{Le constat amiable} is the form drivers fill in together after an accident.

\begin{cousinweb}
\emph{constater} is built on Latin \emph{constat} 'it is certain', from \emph{constāre} 'to stand firm', from \emph{con-} and \emph{stāre} 'to stand' (EtymDB).
\end{cousinweb}

\begin{grammarlens}[title={What you've built}]
One more word for telling, explaining and describing at B1.
\end{grammarlens}

\subsection*{Your turn}
\begin{itemize}
\raggedright
  \item \emph{Say it:} \emph{constater}
  \item \emph{Say it:} \emph{constater}, once more
  \item \emph{Say it:} a sentence of your own with \emph{constater}
  \item \emph{Recall:} say the French for the pedestrian, then the French for the cyclist, then say \emph{constater} again
\end{itemize}

\begin{tcolorbox}[breakable,colback=teal!4,colframe=teal!35!black,title={Before you move on}]
What does constater mean? (\textquotedblleft{}To note, to record (what happened)\textquotedblright{}.) Say it once more.
\end{tcolorbox}
